%
% Bachelor Thesis
% Sven Hodapp
%


% -----------
% 1. Präambel
% -----------


% Allgemeine Einstellungen
% ------------------------
\documentclass[
	pdftex,%              PDFTex verwenden da wir ausschliesslich ein PDF erzeugen.
	a4paper,%             Wir verwenden A4 Papier.
	oneside,%             Einseitiger Druck.
	12pt,%                Grosse Schrift, besser geeignet für A4.
	halfparskip,%         Halbe Zeile Abstand zwischen Absätzen.
	%chapterprefix,%       Kapitel mit 'Kapitel' anschreiben.
	headsepline,%         Linie nach Kopfzeile.
	footsepline,%         Linie vor Fusszeile.
	bibtotocnumbered,%    Literaturverzeichnis im Inhaltsverzeichnis nummeriert einfügen.
	idxtotoc%             Index ins Inhaltsverzeichnis einfügen.
]{report}

\usepackage[utf8]{inputenc}
\usepackage[german]{babel}   % deutsche Silbentrennung
\selectlanguage{german}   % damit Table Of Contents Inhaltsverzeichnis genannt wird

\usepackage{geometry}   % Seitenränder einstellbar
\usepackage{textcomp}   % Sonderzeichen, wie Eurosymbol



% Bilder, Farben, farbige Tabellen
% --------------------------------
\usepackage{graphicx, color, colortbl}
\usepackage{array}       % Erweiterte Tabelleneigenschaften.
%\usepackage{floatflt}   % Bild kann von Text umflossen werden.



% Palatino Schrift
% ----------------
%\usepackage[T1]{fontenc}
%\usepackage[osf]{mathpazo}   % osf aktiviert Mediävalziffern/Minuskelziffern



% Syntax-Highlighting
% -------------------
\definecolor{darkgreen}{rgb}{0,0.5,0}  % needs usepackage color
\usepackage{listings}
\lstset{
	language=Python, 
	basicstyle=\footnotesize\ttfamily, 
	frame=lines, 
	numbers=left, 
	numberstyle=\tiny, 
	numbersep=5pt, 
	breaklines=true, 
	showstringspaces=false, 
	keywordstyle=\color{blue}, 
	commentstyle=\color{darkgreen}, 
	stringstyle=\color{red}
}



% Sonstige Pakete
% ---------------
%\usepackage{anysize}   % Seitenränder verändern
%\usepackage{setspace}   % 1.5em Zeilenabstand \begin{onehalfspacing}
\usepackage{bibgerm}   % Anzeigestil des Literaturverzeichnis (gerabbrv)
\usepackage{paralist}  % Individualisierte Aufzählungen


% Definition von globalen Konstanten
% ----------------------------------
\newcommand{\thema}{
	Vergleich von interner und externer DSL-Technologie zur Entwicklung
	eines Textsatzsystems zur automatischen Dokumentengenerierung.}
\newcommand{\schlagworte}{DSL, Dokumentengenerierung}
\newcommand{\zusammenfassung}{
	Im Zuge dieser Bachelor-Thesis wird ein zeitgemäßer Ansatz für ein Texsatzsystem entwickelt und auf die Probe gestellt, als Vorbild dient TeX.

	Kerntechnologien sind HTML5, CSS3 und Javascript zur Gestaltung und Darstellung eines Dokuments, das dabei entstehende Framework ermöglicht semi-automatische Dokumentengenerierung. Zur Generierung der eigentlichen Dokumenteninhalte werden zwei verschiedene DSL Ansätze Verglichen, zum einen eine interne Scala-DSL, zum anderen eine externe Java-DSL (Xtext).

	Vorteile gegenüber TeX sollen sein: Flexibles Dokument-Templating und interaktive Dokument-Elemente, ermöglicht durch die ubiquitäre HTML-Technologie. Ein moderneres und mächtigeres Sprachkonzept als es LaTeX bietet, mit allen Finessen einer kompletten und verbreiteten Programmiersprache.}
\newcommand{\ausgabedatum}{16.10.2012}
\newcommand{\abgabedatum}{16.01.2013}
\newcommand{\autor}{Sven Hodapp}
\newcommand{\autorStrasse}{Hohentwielstraße 2}
\newcommand{\autorPLZ}{78247 }
\newcommand{\autorOrt}{Hilzingen}
\newcommand{\autorGeburtsort}{Singen/Hohentwiel}
\newcommand{\autorGeburtsdatum}{16.09.1987 }
\newcommand{\prueferA}{Prof. Dr. Marko Boger}
\newcommand{\prueferB}{Dr. Marc Zimmermann}
\newcommand{\firma}{Fraunhofer SCAI}
\newcommand{\studiengang}{Software-Engineering}




% PDF Eigenschaften
% -----------------
\usepackage
[
	colorlinks=false,
	bookmarks = true,
	pdftitle={\thema},
	pdfauthor={\autor},
	pdfsubject={Bachelor Thesis},
	pdfkeywords={\schlagworte},
	urlcolor=blue,
	pdfstartview=FitH
]{hyperref}





% --------------------
% 2. Dokumenten Anfang
% --------------------

\begin{document}



% Deckblatt
% ---------

\begin{titlepage}

\vspace*{-3.5cm}

\begin{flushleft}
\hspace*{-1cm} \includegraphics[width=15.7cm]{titlepages/htwg-logo}
\end{flushleft}

\vspace{2.5cm}

\begin{center}
	\huge{
		\textbf{\thema} \\[4cm]
	}
%\normalsize{\textbf{(Working Draft)} \\}
	\Large{
		\textbf{\autor}} \\[4cm]
	\large{
		\textbf{Konstanz, \abgabedatum} \\[1cm]
	}
	\Huge{
		\textbf{{\sf BACHELORARBEIT}}
	}
\end{center}

\end{titlepage}

\setcounter{page}{2}
\thispagestyle{empty}
{
\setlength{\parskip}{0.5cm}
        \begin{center}
        \textbf{\huge BACHELORARBEIT}

        \textbf{zur Erlangung des akademischen Grades}

        \textbf{\Large Bachelor of Science (B. Sc.)}

        \textbf{an der}

        \textsf{\huge Hochschule Konstanz}\\
        {\small Technik, Wirtschaft und Gestaltung}

        \textsf{\Large Fakult"at Informatik} \\
        Studiengang \studiengang
        \end{center}
}
\begin{center}

\vspace*{2cm}

\begin{tabular}{p{3cm}p{10cm}}
Thema: & \textbf{\large \thema} \\[15ex]
Bachelorkandidat: & \autor, \autorStrasse, \autorPLZ  \autorOrt \\[15ex]
1. Pr"ufer: & \prueferA \\
2. Pr"ufer: & \prueferB \\[25ex]
Ausgabedatum: & \ausgabedatum \\
Abgabedatum: & \abgabedatum \\
\end{tabular}
\end{center}


\begin{center}
{\Large \textbf{Zusammenfassung (Abstract)}}
\end{center}

\bigskip

\begin{center}
	\begin{tabular}{p{2.8cm}p{10cm}}
		Thema: & \thema \\
		 & \\
		Bachelorkandidat: & \autor \\
		 & \\
		Firma: & \firma \\
		 & \\
		Betreuer: & \prueferA  \\[.5ex]
		 &  \prueferB \\
		 & \\
		Abgabedatum: & \abgabedatum \\
		 & \\
		Schlagworte: & \schlagworte \\
		 & \\
	\end{tabular}
\end{center}

\bigskip

\noindent
\zusammenfassung
\chapter*{Ehrenw"ortliche Erkl"arung}
\addcontentsline{toc}{chapter}{Ehrenw"ortliche Erkl"arung}

Hiermit erkl"are ich
\textit{\autor, geboren am \autorGeburtsdatum in \autorGeburtsort}, dass ich\\

\begin{tabular}{lp{12cm}}
(1) & meine Bachelorarbeit mit dem Titel \\[1em]
& \textbf{\thema} \\[1em]
& bei der \firma\ unter Anleitung von \prueferA\ selbst"andig und ohne fremde Hilfe angefertigt und keine anderen als die angef"uhrten Hilfen benutzt habe;\\[1em]
(2) & die "Ubernahme w"ortlicher Zitate, von Tabellen, Zeichnungen, Bildern und
Programmen aus der Literatur oder anderen Quellen (Internet) sowie die Verwendung
der Gedanken anderer Autoren an den entsprechenden Stellen innerhalb der Arbeit
gekennzeichnet habe.\\
\end{tabular}

\vspace*{1cm}

\noindent
Ich bin mir bewusst, dass eine falsche Erkl"arung rechtliche Folgen haben wird.\\

\vspace*{3cm}

\noindent
Konstanz, \abgabedatum \hfill \begin{tabular}{c} \\ \\ \rule{5cm}{1pt} \\ (Unterschrift)\end{tabular}




% Inhaltsverzeichnis anzeigen
% ---------------------------
\tableofcontents
%\listoffigures
%\listoftables




% ---------
% 3. Inhalt
% ---------

\chapter{Vorwort}

Vorwort.\cite{ref}


\addcontentsline{toc}{section}{Literaturverzeichnis}
\bibliographystyle{geralpha}
\bibliography{bib/references}

\end{document}


